%===============================================================================
% Experience Report
%===============================================================================
% $Id$
%===============================================================================


%===============================================================================
% Configuration
%===============================================================================


%-------------------------------------------------------------------------------
% \documentclass and \usepackage directives
%-------------------------------------------------------------------------------
\documentclass[a4paper,fleqn,titlepage]{article} 
%\usepackage{ngerman}
\usepackage[latin1]{inputenc}
\usepackage[T1]{fontenc}
\usepackage[small,hang,bf]{caption2}
\usepackage{fancyhdr}
\usepackage[nice]{nicefrac}
\usepackage{color,listings}
\usepackage{alltt}


% Compilation with latex or pdflatex?
\newif\ifpdf 
\ifx\pdfoutput\undefined 
  \pdffalse
\else
  \pdfoutput=1 
  \pdftrue 
\fi 

% Compilation with pdflatex
\ifpdf
 
  \usepackage[pdftex]{graphicx}

  \usepackage[
    pdftex,
    a4paper,
    bookmarks,
    pdfstartview=FitH,    % starts with page width
    bookmarksopen,        % opens index
    bookmarksnumbered,    % index with numbering
    colorlinks,           % links with color, otherwise with border
    linkcolor=blue,       % Standard red
    citecolor=blue,       % Standard green
    urlcolor=magenta,     % Standard cyan
    filecolor=blue
  ]{hyperref} 

  \pdfinfo{
    /Title      (Experience Report)
    /Author     (Thomas Weibel, Martin Luder, Michael K�ser)
    /Subject    (Eiffel programming)
    /Keywords   (Programming, EiffelRSS)
  }

  % Use default Acrobat reader fonts
  \usepackage{mathpazo}

  % Use CM fonts (increases document size)
  % \usepackage{ae}

% Compilation with latex
\else 

  \usepackage{graphicx} 

\fi


%-------------------------------------------------------------------------------
% Configure \maketitle
%-------------------------------------------------------------------------------
\title{EiffelRSS \\ Experience Report}
\author{
  Michael K\"aser <kaeserm@student.ethz.ch>
  \and 
  Martin Luder <luderm@student.ethz.ch>
  \and 
  Thomas Weibel <weibelt@student.ethz.ch>
}
\date{\today}


%-------------------------------------------------------------------------------
% Configure fancyhdr
%-------------------------------------------------------------------------------
\pagestyle{fancy}

\renewcommand{\headrulewidth}{0.1 pt}
\renewcommand{\footrulewidth}{0.1 pt}

\fancypagestyle{plain}{
  \lhead{EiffelRSS}
  \chead{}
  \rhead{\thepage}
  \lfoot{EiffelRSS}
  \cfoot{}
  \rfoot{Experience Report}
}

\lhead{EiffelRSS}
\chead{}
\rhead{\thepage}

\lfoot{EiffelRSS}
\cfoot{}
\rfoot{Experience Report}


%-------------------------------------------------------------------------------
% Configure listings
%-------------------------------------------------------------------------------
\lstset{showstringspaces=false,
  breaklines=true,
  breakindent=0pt,
  prebreak=\mbox{\tiny$\searrow$},
  postbreak=\mbox{{\color{blue}\tiny$\rightarrow$}},
  frame=trBL,
  framerule=0.75pt,
  framesep=4pt,
  rulesep=0.75pt
}


%-------------------------------------------------------------------------------
% Common configuration
%-------------------------------------------------------------------------------
\setlength{\parindent}{0em}
\setlength{\parskip}{1.5ex plus0.5ex minus0.5ex}
\sloppy
\setlength{\mathindent}{0em}


%-------------------------------------------------------------------------------
% Commandos
%-------------------------------------------------------------------------------
\newcommand{\hr}{\rule{\textwidth}{1pt}}


%===============================================================================
% Document
%===============================================================================
\begin{document}

\begin{titlepage}
  \newlength{\centeroffset}
  \setlength{\centeroffset}{-0.5\oddsidemargin}
  \addtolength{\centeroffset}{0.5\evensidemargin}

  \thispagestyle{empty}

  \noindent\includegraphics[width=\textwidth]{../figures/big_ETH}\\[-3mm]
  \hr

  \vspace*{\stretch{1}}

  \makebox[0pt][l]{
    \begin{minipage}{\textwidth}
      \flushright{
        \Huge\bfseries EiffelRSS
      }

      \noindent\rule{\textwidth}{3pt}\\[2.5ex]

      \hfill\emph{
        \Large Experience Report
      }
    \end{minipage}
  }

  \vspace{\stretch{1}}

  \makebox[0pt][l]{
    \begin{minipage}{\textwidth}
      \flushright{
        \bfseries 
        Michael K\"aser <kaeserm@student.ethz.ch>\\[0.3ex]
        Martin Luder <luderm@student.ethz.ch>\\[0.3ex]
        Thomas Weibel <weibelt@student.ethz.ch>\\[0.3ex]
      }
    \end{minipage}
  }

  \vspace{\stretch{1}}

  \noindent\hr\\[1mm]
  \includegraphics[width=\textwidth]{../figures/big_inf}
\end{titlepage}

\begin{abstract}
  The experience report describes the teamwork, things learned,
  difficulties, etc. from each project members' point of view.
\end{abstract}

\section{Michael K�ser}
\label{sec:michael}

I learned quite a lot of new things during the work, mainly in terms
of teamwork. I haven't really programmed in a team before this project
and always thought that it must be quite troublesome, but Thomas had a
lot of experience and knew some very good tools which made everything
very easy, for example the Wiki, Subversion and his great \LaTeX{}
templates. We also could divide the work quite good in three parts.

This project was also the first time that I actually thought a lot
about the design of my classes before writing them. This wasn't a
problem for me, but I had a lot of difficulties with the EiffelNet
library, because it is so poorly documented and implemented. Without
any examples you have to read all the code to understand how it
actually works.

In the end I'm really happy with the result of this project. The
design is fine and the library is useful. The project was a good
experience and definitively the most interesting thing at ETH until
now.


\section{Martin Luder}
\label{sec:martin}

During the work on this project, I've learned a lot: Newsreader was my
first GUI that really does something. I've never worked with
Subversion (or any similar tool) before. Because of Subversion, the
teamwork worked flawless, and when there were any questions about the
library, I'd just post it to the Wiki and would get an answer quickly.

I didn't really have many difficulties, because the biggest part of
the work was with EiffelVision, which is quite simple to handle
(although it didn't support all the things I wanted to do, like GIF
support). But some issues in EiffelNet and other parts of the Eiffel
library annoyed me nonetheless, because Newsreader just wouldn't load
some feeds and it took some time until Michael and Thomas worked
around them.

Overall, this project was the best thing during the whole term. It was
a lot of fun actually doing something concrete after all the theory!


\section{Thomas Weibel}
\label{sec:thomas}

To work on this project was really exciting and I learned a lot,
especially about Eiffel programming, teamwork and syndication
standards or better the lack of standards.

The teamwork was excellent. We discussed most design issues in the
Wiki and had a weekly project meeting. The Wiki helped a lot to
coordinate our work. To manage our source code, we used Subversion
which is an extremely good tool for such a quick evolving project like
EiffelRSS. Subversion allowed us to easily restructure our clusters
and classes without losing the whole revision history.

Most problems I encountered had to do with ISE Eiffel. I really
enjoyed programming with Eiffel but the Eiffel library of ISE is very
cumbersome:

\begin{itemize}
\item No real stream concept, e.g. input and output are no files
\item No sortable data structures to implement sorting by different
  criterias, e.g. sort a list of addresses according to requirements
  by name, phone or city.
\item Only very meagre network support, e.g. no HTTP 1.1 client/server
  implementation
\end{itemize}

Fortunately ISE delivers Gobo with EiffelStudio.

In conclusion, I think we have every reason to be proud of what we
have achieved especially by considering the limited time we had
available.

\end{document}
